\documentclass[11pt]{article}
\usepackage{amsmath,amssymb,amsthm}
\newtheorem{lemma}{Lemma}
\title{A Fixed-Pool Coalescing-Query Audit}
\author{}
\date{}
\begin{document}
\maketitle

\paragraph{Scope.}
This short written lemma excludes one compact-limit shortcut. It does
not exclude scale-dependent poles or residues, does not bound the number
of patterns for arbitrary fixed query intervals, and does not establish
either requested HeadComplexity asymptotic result. It is not Lean checked.

\begin{lemma}
Let $\mathcal F$ be a finite family of real-analytic functions in a
neighborhood of $T$. Fix nonzero real numbers $u_1,\ldots,u_q$.
For all sufficiently small $\eta>0$, the strict sign vectors
\[
 \bigl(\operatorname{sgn}f(T+\eta u_j)\bigr)_{j=1}^q,
 \qquad f\in\mathcal F,
\]
belong to a set of at most four patterns: the two constant patterns and
the two patterns $\pm(\operatorname{sgn}u_j)_j$.
\end{lemma}
\begin{proof}
An identically zero function supplies no strict sign vector. Otherwise,
let $k$ be the order of its first nonzero Taylor coefficient $a$ at $T$.
For each fixed nonzero $u_j$,
\[
 f(T+\eta u_j)=a\eta^ku_j^k+o(\eta^k).
\]
The leading coefficient is nonzero. Thus for sufficiently small positive
$\eta$ its sign is $\operatorname{sgn}(a)\operatorname{sgn}(u_j)^k$.
Take a common positive threshold over the finitely many functions and
queries. Even $k$ gives a constant pattern; odd $k$ gives the signed
$u$ pattern.
\end{proof}

In particular, fix poles $p_h<T$ and real weights $a_h$. Applying the
lemma to all $2^M$ functions
$f_s(t)=\sum_{h=1}^M s_h a_h/(t-p_h)$ shows that a fixed weighted
Cauchy sign cube has at most four strict output patterns as the query
interval contracts to $T$. Letting the pool itself depend on the query
scale invalidates this application, and is permitted in the research
problem. Likewise, sending poles to zero while holding a nonzero query
interval fixed is a different limit: its leading coefficient can be a
nonconstant rational function of the query time.

\paragraph{Exact finite diagnostic.}
The accompanying script \texttt{compact-cauchy-signcube-probe.py}
exhausts all $2^M$ sign choices for $M=4,6,8,10,12,14,16$.
It uses rational equally spaced poles and query times, and six displayed
weight choices: equal weights; $T-p_h$; a small perturbation thereof;
barycentric weights; and two barycentric weights multiplied by prescribed
polynomial values. Every denominator is positive. Each rational row is
cleared to integers, and Gray-code updates evaluate signs with exact
integer arithmetic. In these configurations the counts are always two
or four. This is computational evidence for those explicit families,
not a bound over all scale-dependent pools.

\paragraph{Entropy budget warning.}
If a proposed recursive decoder gives $\exp(cH\log H)$ distinct
patterns on $O(H\log H)$ queries, Sauer's lemma does imply a shattered
subset of size $\Omega(H\log H)$. However,
$\exp(cH\log\log H)$ patterns on the same query count do not by
themselves imply a superconstant ratio of shattered size to $H$.
Indeed a class of VC dimension $d=\Theta(H)$ can already have up to
\[
 (eq/d)^d=\exp(\Theta(H\log\log H)),
 \qquad q=\Theta(H\log H).
\]
A weaker entropy recurrence must therefore also reduce the retained query
count, or prove a further distance or projection property. Merely counting
permutations or distinct scalar states does not supply this missing step.
\end{document}
